\documentclass[11pt,a4paper]{article}

\usepackage[utf8]{inputenc}
\usepackage[T1]{fontenc}
\usepackage{lmodern}
\usepackage[slovak]{babel}
\usepackage[a4paper,margin=2.2cm]{geometry}
\usepackage{booktabs}
\usepackage{tabularx}
\usepackage{array}
\usepackage{longtable}
\usepackage{enumitem}
\usepackage[hidelinks]{hyperref}
\usepackage[protrusion=true,expansion=false]{microtype}
\usepackage{parskip}
\usepackage{cite}

% Spoločné údaje projektu — pri zmene upravte iba tento súbor.
\newcommand{\projectTitleSK}{Detekcia falošných online recenzií na základe jazykových a behaviorálnych príznakov}
\newcommand{\projectTitleEN}{Detection of Fake Online Reviews Based on Linguistic and Behavioral Features}
\newcommand{\projectAcronym}{FAKERADAR}
\newcommand{\projectId}{Zatiaľ nepridelené (študentský návrh)}
\newcommand{\projectDate}{01.11.2026}
\newcommand{\projectStart}{01/2027}
\newcommand{\projectEnd}{12/2029}

\newcommand{\applicant}{Slovenská technická univerzita v Bratislave, Fakulta informatiky a informačných technológií}
\newcommand{\applicantShort}{FIIT STU}
\newcommand{\applicantAddress}{Ilkovičova 6276/2, 842 16 Bratislava 4}
\newcommand{\applicantICO}{00397687}
\newcommand{\applicantDIC}{2020845255}
\newcommand{\applicantICDPH}{SK2020845255}
\newcommand{\applicantEmail}{info@fiit.stuba.sk}
\newcommand{\applicantPhone}{+421 2 210 44 111}
\newcommand{\statutory}{prof. Ing. Ivan Kotuliak, PhD.}
\newcommand{\statutoryRole}{dekan}
\newcommand{\piName}{Ing. Ivan Kapustík}
\newcommand{\piEmail}{ivan.kapustik@stuba.sk} % e-mail potvrdený v oficiálnych materiáloch FIIT

\newcommand{\budgetAPVV}{268\,000}
\newcommand{\budgetCoFin}{32\,000}
\newcommand{\budgetTotal}{300\,000}

% Anotácia = presné znenie schváleného zamerania (odovzdané 04.10.2026, hodnotené cvičiacim 5/5 b);
% zdroj: old_zameranie_clanok.tex (abstrakt).
\newcommand{\projectAnnotation}{%
Online recenzie výrazne ovplyvňujú rozhodovanie spotrebiteľov, no ich dôveryhodnosť narúšajú falošné recenzie generované firmami, platenými recenzentmi či automatizovanými nástrojmi. Tento článok prináša prehľad súčasných metód detekcie falošných recenzií so zameraním na porovnanie jazykových (textových) a behaviorálnych (na správaní recenzenta založených) príznakov.}


\setlength{\parindent}{0pt}
\setlength{\emergencystretch}{3em}
\newcommand{\raggedcell}{\raggedright\arraybackslash}

\begin{document}
\sloppy

\begin{titlepage}
\centering
\vspace*{1.2cm}
{\Large\bfseries VV-F — Vecný zámer projektu}\\[0.3em]
{\large Záväzná osnova pre aplikovaný výskum a vývoj}\\[2cm]
{\LARGE\bfseries \projectTitleSK}\\[0.8em]
{\large \projectTitleEN}\\[1.5cm]
\begin{tabular}{@{}>{\raggedright\arraybackslash}p{4.2cm} p{9.5cm}@{}}
\textbf{Akronym:} & \projectAcronym\\
\textbf{Evidenčné číslo:} & \projectId\\
\textbf{Žiadateľ:} & \applicant\\
\textbf{Zodpovedný riešiteľ:} & \piName\\
\textbf{Dátum:} & \projectDate\\
\textbf{Obdobie riešenia:} & \projectStart\ -- \projectEnd\\
\end{tabular}
\vfill
{\footnotesize
Dokument je modelovým návrhom v rozsahu približne 50--60\,\% očakávanej finálnej verzie. Personálne a ekonomické údaje označené ako „doplniť“ treba pred finálnym podaním aktualizovať.
}
\end{titlepage}

\section{Excelentnosť}

\subsection{Projektový zámer a aplikačná úroveň}

Online recenzie sú jedným z najvplyvnejších informačných kanálov v elektronickom obchode. Spotrebitelia na ich základe porovnávajú produkty, predajcov a ceny, podniky zase budujú reputáciu a konkurenčnú výhodu. Táto ekonomická hodnota vytvára motiváciu manipulovať recenzným priestorom prostredníctvom platených recenzentov, konkurenčných útokov, koordinovaných skupín účtov alebo automatizovaných nástrojov. Falošné recenzie preto poškodzujú spotrebiteľov, znižujú dôveryhodnosť platforiem a narúšajú férovú súťaž.

Projekt reaguje na medzeru medzi akademickými detektormi a potrebami praxe. Existujúce riešenia sú prevažne anglocentrické, založené na povrchových jazykových príznakoch a v podmienkach generatívnych modelov strácajú spoľahlivosť. Prevádzkovatelia platforiem zároveň potrebujú nielen presný, ale aj vysvetliteľný nástroj, ktorý dokáže rozhodnutie o označení recenzie zdôvodniť a ktorý je možné prevádzkovať bez neprimeraného zásahu do súkromia používateľov.

Projekt preto vyvinie hybridný detektor falošných recenzií, ktorý spája jazykové, behaviorálne a grafové príznaky, doplní ho o vrstvu vysvetlenia a overí v pilotnej prevádzke. Výstupom bude otvorený slovensko-anglický anotovaný dataset, metodika anotácie a hodnotenia, referenčný model, prototyp API a správa pre prax. Projekt začína na úrovni TRL~3 a končí na úrovni TRL~6: od overených laboratórnych postupov k funkčnému prototypu demonštrovanému v relevantnom prostredí.

\subsection{Aktuálnosť a originalita}

Aktuálnosť témy je daná tromi súbežnými trendmi. Prvým je rast elektronického obchodu a závislosti spotrebiteľov od recenzií. Druhým je dostupnosť veľkých jazykových modelov, ktoré dokážu generovať plynulé a štylisticky presvedčivé recenzie v takmer neobmedzenom množstve \cite{zhao2025}. Tretím je regulácia: európsky rámec zakazuje zavádzajúce recenzie a posilňuje transparentnosť platforiem, čím vzniká dopyt po obhájiteľných nástrojoch moderovania.

Originalita projektu spočíva v piatich prvkoch. (i)~Vytvoríme prvý otvorený slovensko-anglický benchmark pre detekciu falošných recenzií s viacnásobnou anotáciou. (ii)~Spojíme textové, behaviorálne a grafové príznaky v jednom modeli a systematicky porovnáme ich prínos. (iii)~Budeme explicitne modelovať recenzie generované LLM, nielen ľuďmi písaný spam. (iv)~Vysvetliteľnosť bude rovnocenným cieľom, nie dodatočnou funkciou. (v)~Výstupy budú otvorené a reprodukovateľné, aby ich mohli využiť ďalšie tímy aj prevádzkovatelia platforiem.

Nejde teda o ďalší izolovaný klasifikátor, ale o metodiku a nástroj, ktorý prepája výskum s moderovaním v praxi. Práve kombinácia dát, hybridného modelovania a vysvetliteľnosti vytvára priestor na originálny príspevok v slovenskom aj medzinárodnom kontexte.

\subsection{Stav poznania a vedecké základy}

Výskum detekcie falošných recenzií sa rozvíja od práce Jindala a Liu~\cite{jindal2008}, ktorí vymedzili taxonómiu recenzného spamu a ukázali, že duplicitné texty sú silným indikátorom manipulácie. Ott a kol.~\cite{ott2011} vytvorili prvý štandardizovaný dataset klamlivých recenzií a demonštrovali, že n-gramy a psycholingvistické lexikóny prekonávajú ľudských posudzovateľov. Ich práca zároveň upozornila na metodologický problém: pri autenticky pôsobiacich textoch sa klasifikátor môže učiť idiosynkrázie konkrétnych autorov namiesto všeobecného pojmu falošnosti.

Ekonomické motivácie analyzovali Luca a Zervas~\cite{luca2016}, ktorí dokumentovali vzťah medzi reputáciou podniku a objednávaním falošných pozitívnych recenzií. Mukherjee a kol.~\cite{mukherjee2013} ukázali, že behaviorálne príznaky môžu byť informatívnejšie než samotný text. Samostatnú líniu predstavujú práce o koordinovaných skupinách a recenzných sieťach: Mukherjee a kol.~\cite{mukherjee2012} odhaľovali falošné recenzentské skupiny, Rayana a Akoglu~\cite{rayana2015} spojili recenzie, recenzentov a produkty do jednej siete a využili šírenie podozrenia medzi prepojenými entitami.

Prehľadové práce \cite{crawford2015} a \cite{heydari2015} zjednotili terminológiu a rozdelili prístupy podľa typu príznakov a úrovne rozhodovania. Barbado a kol.~\cite{barbado2019} overili aplikačný rozmer na doméne spotrebnej elektroniky a Li a kol.~\cite{li2014} ukázali, že existujú relatívne univerzálne, ľudsky interpretovateľné pravidlá prekračujúce hranice datasetu. Zhao a kol.~\cite{zhao2025} napokon dokumentovali, že AI-generované recenzie sú plynulejšie a štylisticky konzistentnejšie než ľudské, čo oslabuje čisto textové detektory.

Napriek tomuto pokroku zostávajú medzery, na ktoré projekt reaguje. Chýba otvorený slovensko-anglický dataset, systematické porovnanie jazykových a behaviorálnych príznakov v podmienkach LLM, prakticky použiteľný hybridný detektor s vysvetliteľnými výstupmi a metodika hodnotenia, ktorá zohľadňuje časový a doménový posun. Vedecké základy projektu preto tvoria štyri oblasti: spracovanie prirodzeného jazyka a štylometria, behaviorálna analýza a dolovanie v grafoch, strojové učenie s dôrazom na kalibráciu a generalizáciu a nakoniec výskum vysvetliteľnosti a používateľskej dôvery.

\subsection{Ciele projektu}

Ciele sú formulované merateľne a viažu sa na konkrétne výstupy a termíny.

\begin{table}[h!]
\centering
\small
\begin{tabularx}{\textwidth}{@{}p{0.8cm} X p{4.2cm} p{1.5cm}@{}}
\toprule
\# & Cieľ & Ukazovateľ & Termín \\
\midrule
O1 & Vytvoriť otvorený slovensko-anglický anotovaný dataset & minimálne 15\,000 recenzií; $\kappa \geq 0{,}7$ & M12 \\
O2 & Vyvinúť jazykový model & F1 $\geq 0{,}80$ v rámci domény & M18 \\
O3 & Vyvinúť behaviorálny a grafový model & F1 $\geq 0{,}75$; odhalenie skupín & M24 \\
O4 & Skonštruovať hybridný model & F1 $\geq 0{,}85$ in-domain; $\geq 0{,}75$ cross-domain & M30 \\
O5 & Overiť prototyp v praxi & pilot s 1--2 platformami; vysvetlenie top-$k$ príznakov & M36 \\
\bottomrule
\end{tabularx}
\end{table}

Ciele sú realistické vzhľadom na rozsah projektu, dostupné datasety a skúsenosti tímu. Zároveň sú dostatočne ambiciózne: vyžadujú nielen trénovanie modelov, ale aj tvorbu dát, anotáciu, evaluáciu a pilotné overenie. Splnenie cieľov budeme priebežne vyhodnocovať na míľnikoch M12, M24 a M36.

\subsection{Metodológia}

Metodológia je navrhnutá ako sekvencia ôsmich krokov, ktoré sa čiastočne prekrývajú.

\textbf{1. Zber dát.} Využijeme verejné datasety (napr.\ Ott a kol.~\cite{ott2011}), verejne dostupné recenzie z platforiem a dátové zdroje partnerských organizácií. Doplníme ich syntetickými recenziami generovanými viacerými LLM, aby sme pokryli nový typ hrozby. Pri všetkých dátach budeme dodržiavať licenčné podmienky a GDPR.

\textbf{2. Anotácia.} Vytvoríme anotačné guidelines s definíciami typov falošných recenzií, hraničných prípadov a neurčitých príkladov. Každú recenziu budú anotovať aspoň dvaja nezávislí anotátori; nezhody vyrieši tretí anotátor. Kvalitu budeme merať Cohenovou kappa a pravidelnými kalibračnými stretnutiami.

\textbf{3. Predspracovanie.} Text normalizujeme, identifikujeme jazyk, odstránime duplicity a citlivé údaje. Pre slovenčinu využijeme dostupné nástroje na tokenizáciu, lemmatizáciu a morfologickú analýzu; pri angličtine použijeme štandardné NLP nástroje.

\textbf{4. Extrakcia príznakov.} Jazykové príznaky zahŕňajú lexikálnu bohatosť, distribúciu slovných druhov, sentiment, subjektivitu, štylometrické markery a transformer-ové reprezentácie. Behaviorálne príznaky zahŕňajú históriu recenzenta, frekvenciu, burstiness, ratingovú odchýlku a overený nákup. Grafové príznaky zahŕňajú bipartitný graf recenzent--produkt, ko-recenzné skupiny a miery centrality.

\textbf{5. Modelovanie.} Ako baselines použijeme logistickú regresiu, SVM a gradientové boosting. Následne natrénujeme transformer (BERT alebo jeho viacjazyčný variant) a grafovú neurónovú sieť. Hybridný model vznikne early fusion (spoločný vektor príznakov) a late fusion (vážená kombinácia rozhodnutí); porovnáme oba varianty.

\textbf{6. Evaluácia.} Použijeme stratifikovanú krížovú validáciu, časové rozdelenie a rozdelenie podľa domén. Sledujeme precision, recall, F1, AUROC, mieru falošných pozitív pri fixnom zachytení a kalibráciu (ECE). Súčasťou je aj testovanie na novom jazyku a na adversariálnych príkladoch.

\textbf{7. Vysvetliteľnosť.} Pre globálne vysvetlenie použijeme SHAP a permutation importance, pre lokálne vysvetlenie LIME a attention. Overíme, či vysvetlenia zodpovedajú intuícii moderátorov a či zvyšujú ich dôveru v model.

\textbf{8. Pilot.} Prototyp nasadíme v režime tieňového rozhodovania, aby neovplyvňoval reálne rozhodnutia. Zbierame spätnú väzbu moderátorov, meriame čas rozhodovania a mieru falošných pozitív. Výsledky použijeme na finálnu kalibráciu a na správu pre prax.

Etický rámec zahŕňa minimalizáciu osobných údajov, anonymizáciu, informovaný súhlas pri partnerských dátach a nezávislý pohľad na riziko zneužitia detektora na cenzúru legitímnych recenzií.

\subsection{Predchádzajúce výsledky a kompetencie}

Zodpovedný riešiteľ \piName\ pôsobí na FIIT STU (Ústav informatiky a softvérového inžinierstva) a dlhodobo sa venuje softvérovému inžinierstvu, práci s dátami a vedeniu študentských projektových tímov (predmet Tímový projekt). Expertízu tímu v oblasti spracovania prirodzeného jazyka a strojového učenia zabezpečí spolupráca so špecialistami fakulty a mladými riešiteľmi. Uvedieme najviac 5 najvýznamnejších výstupov za posledných 5 rokov a najviac 3 projekty, ktoré zodpovedný riešiteľ viedol. Položky sú zoradené od najvýznamnejšej po menej významnú (podľa vyjadrenia Ing. Kapustíka); konkrétne údaje doplníme pred podaním podľa životopisu zodpovedného riešiteľa.

Jeho víziou je budovať výskumný tím zameraný na dôveryhodnú analýzu textu a správania, systematicky rozvíjať expertízu v oblasti vysvetliteľného strojového učenia a vytvárať priestor pre kariérny rast mladých výskumníkov — od semestrálnych projektov cez diplomové práce až po doktorandské štúdium a vlastné grantové projekty.

\begin{table}[h!]
\centering
\small
\begin{tabularx}{\textwidth}{@{}p{0.8cm} X p{2.0cm}@{}}
\toprule
\# & Najvýznamnejší výstup za posledných 5 rokov (max. 5, od najvýznamnejšieho) & Rok \\
\midrule
V1 & \emph{doplniť názov publikácie/výstupu a jeho význam} & doplniť \\
V2 & \emph{doplniť názov publikácie/výstupu a jeho význam} & doplniť \\
V3 & \emph{doplniť názov publikácie/výstupu a jeho význam} & doplniť \\
V4 & \emph{doplniť názov publikácie/výstupu a jeho význam} & doplniť \\
V5 & \emph{doplniť názov publikácie/výstupu a jeho význam} & doplniť \\
\bottomrule
\end{tabularx}
\end{table}

\begin{table}[h!]
\centering
\small
\begin{tabularx}{\textwidth}{@{}p{0.8cm} X p{2.2cm} p{1.6cm}@{}}
\toprule
\# & Projekt (max. 3, najvýznamnejšie; názov, schéma, roky, rozpočet) & Rola (viedol) & Význam \\
\midrule
P1 & \emph{doplniť} & \emph{doplniť} & \emph{doplniť} \\
P2 & \emph{doplniť} & \emph{doplniť} & \emph{doplniť} \\
P3 & \emph{doplniť} & \emph{doplniť} & \emph{doplniť} \\
\bottomrule
\end{tabularx}
\end{table}

Organizácia disponuje laboratóriami, výpočtovou infraštruktúrou a prístupom k vedeckým databázam. Fakulta má skúsenosti s riešením grantových projektov, s vedením bakalárskych a diplomových prác a s publikovaním v oblasti informatiky. Pre projekt je dôležitá aj medzinárodná spolupráca a možnosť konzultovať výsledky s odborníkmi z praxe.

\subsection{Zapojenie mladých pracovníkov}

Na riešení sa podieľajú štyria študenti druhého stupňa, ktorí tvoria 80\,\% riešiteľského tímu. Matvii Nosenko sa zameria na NLP a modely strojového učenia, Volodymyr Nedoluha na zber a spracovanie dát, Danylo Nikolenko na backend a API a Rostyslav Mykytiuk na evaluáciu a vizualizáciu výsledkov. Expertízu jednotlivých riešiteľov tvoria kvalifikačné a semestrálne práce v oblasti spracovania prirodzeného jazyka, analýzy dát a softvérového inžinierstva; pred podaním ju doplníme o konkrétne výsledky a skúsenosti. Každý z nich bude mať možnosť publikovať ako spoluautor, zúčastniť sa konferencie a získať praktické skúsenosti s riešením aplikovaného projektu. Plán rozvoja zahŕňa pravidelné odborné semináre, revízie kódu a mentoring zo strany zodpovedného riešiteľa.

\section{Dopad}

\subsection{Prínosy pre poznanie, aplikovaný výskum a inovácie}

Projekt prinesie štyri druhy prínosov. Prvým je otvorený slovensko-anglický dataset a benchmark, ktorý umožní reprodukovateľné porovnávanie metód a podporí ďalší výskum v slovenskom kontexte. Druhým je hybridná metodika, ktorá systematicky meria prínos jazykových, behaviorálnych a grafových príznakov a ukazuje, kedy sa ktorá skupina vyplatí. Tretím je model s vysvetliteľnými výstupmi, ktorý prepája presnosť s dôveryhodnosťou pre moderátorov. Štvrtým je prototyp API a dashboard, ktorý demonštruje použiteľnosť v praxi.

Z hľadiska inovácií ide o novú službu pre prevádzkovateľov platforiem: modulárny detektor, ktorý možno nasadiť ako súčasť moderovacieho procesu alebo ako samostatnú analytickú vrstvu. Riešenie nie je viazané na konkrétnu platformu a jeho architektúra umožňuje rozšírenie o ďalšie jazyky a domény.

\subsection{Využiteľnosť výsledkov v SR a zahraničí}

Na Slovensku môžu výsledky využiť prevádzkovatelia e-shopov a trhovísk, recenzné portály, profesijné asociácie a orgány ochrany spotrebiteľa. V zahraničí je potenciál najmä v krajinách EÚ, kde platí spoločný regulačný rámec pre zavádzajúce recenzie a transparentnosť platforiem. Dataset a metodika sú jazykovo rozšíriteľné, takže ich možno adaptovať na český, poľský alebo maďarský jazyk.

Výsledky budú využiteľné v troch režimoch: (i)~offline analýza historických dát pre audit a reporting; (ii)~online asistencia moderátorovi s vysvetlením a prioritizáciou podozrivých recenzií; (iii)~monitorovanie trendov a nových typov útokov. Každý režim má iné požiadavky na latenciu, vysvetliteľnosť a ochranu údajov, čo projekt zohľadní v návrhu API.

\subsection{Ekonomický a spoločenský prínos}

Falošné recenzie spôsobujú priame straty spotrebiteľom, poškodzujú poctivých predajcov a zvyšujú náklady platforiem na moderovanie. Presnejší a vysvetliteľný detektor môže tieto náklady znížiť, zvýšiť dôveru v recenzie a zlepšiť kvalitu rozhodovania. Prínos je zároveň spoločenský: ochrana spotrebiteľov, transparentnosť a férová súťaž.

Kvantifikácia prínosu bude predmetom pilotu. Budeme sledovať počet odhalených podozrivých recenzií, mieru falošných pozitív, čas moderátora na jedno rozhodnutie a subjektívnu dôveru moderátorov. Ekonomický prínos vyjadríme ako odhad úspory času a zníženia počtu sťažností; pôjde o konzervatívny odhad, nie o marketingové číslo.

\subsection{Opatrenia na maximalizáciu dopadu a komunikácia}

Projekt bude priebežne zverejňovať výsledky formou odborných článkov, preprintov, workshopov a krátkych správ pre prax. Dataset a kód sprístupníme za jasných licenčných podmienok. Pre prevádzkovateľov pripravíme praktického sprievodcu nasadením a pre regulátorov stručný policy brief. Komunikácia bude prebiehať cez web fakulty, odborné podujatia a partnerské kanály.

Meranie dopadu bude na troch úrovniach: vedeckej (publikácie, citácie, využitie datasetu), technologickej (počet nasadení, latencia, presnosť) a spoločenskej (počet chránených spotrebiteľov, zníženie sťažností). Tieto ukazovatele budeme vyhodnocovať na konci projektu a priebežne na míľnikoch.

\section{Implementácia}

\subsection{Projektový plán, pracovné balíky a výstupy}

Projekt je rozdelený na sedem pracovných balíkov, ktoré sa čiastočne prekrývajú. Každý balík má zodpovedného riešiteľa, výstupy a míľnik.

\begin{table}[h!]
\centering
\small
\begin{tabularx}{\textwidth}{@{}p{0.8cm} X p{3.0cm} p{1.4cm}@{}}
\toprule
WP & Názov a obsah & Výstup & Mesiace \\
\midrule
WP1 & Manažment, koordinácia, reporting, riadenie rizík & plán, správy & 1--36 \\
WP2 & Zber, anotácia a sprístupnenie dát & dataset, guidelines & 1--12 \\
WP3 & Jazykové modely a baseline & model, evaluácia & 7--18 \\
WP4 & Behaviorálne a grafové modely & model, evaluácia & 13--24 \\
WP5 & Hybridný model a vysvetliteľnosť & prototyp & 19--30 \\
WP6 & Pilotné overenie a evaluácia & pilotná správa & 25--36 \\
WP7 & Diseminácia a komunikácia & publikácie, web & 1--36 \\
\bottomrule
\end{tabularx}
\end{table}

Výstupy sú navrhnuté tak, aby na seba nadväzovali: dataset a guidelines sú predpokladom pre modely, modely pre hybridný prototyp a prototyp pre pilot. Každý výstup má jasného vlastníka a kontrolu kvality.

\subsection{Harmonogram a míľniky}

\begin{table}[h!]
\centering
\small
\begin{tabularx}{\textwidth}{@{}p{1.2cm} X p{1.8cm}@{}}
\toprule
Míľnik & Popis & Termín \\
\midrule
M6 & Schválená metodika anotácie a etický rámec & 06/2027 \\
M12 & Dataset SK/EN s minimálne 15\,000 recenziami & 12/2027 \\
M18 & Jazykový model a prvá publikácia & 06/2028 \\
M24 & Behaviorálny a grafový model & 12/2028 \\
M30 & Hybridný model a prototyp API & 06/2029 \\
M36 & Pilotná správa, finálne vyhodnotenie a druhá publikácia & 12/2029 \\
\bottomrule
\end{tabularx}
\end{table}

\subsection{Projektový manažment}

Projekt bude riadený zodpovedným riešiteľom, ktorý rozhoduje o prioritách, rozpočte a personálnych otázkach. Operatívne riadenie zabezpečí projektový koordinátor z tímu. Budeme používať pravidelné týždenné stretnutia, mesačné revízie a štvrťročné hodnotiace správy. Úlohy budeme evidovať v Jira a kód v Git repozitári; každá zmena bude mať issue key a code review.

Kontrolné mechanizmy zahŕňajú: (i)~kontrolu kvality dát pred uzamknutím datasetu; (ii)~reprodukovateľ\-nosť experimentov cez fixné seeds a verzovanie; (iii)~nezávislú revíziu evaluačného protokolu; (iv)~etickú kontrolu pred prácou s partnerskými dátami. Riziká budeme priebežne aktualizovať v registri rizík a na mesačných revíziách.

\subsection{Riziká implementácie a ich zmiernenie}

\begin{table}[h!]
\centering
\small
\begin{tabularx}{\textwidth}{@{}p{3.3cm} X X@{}}
\toprule
Riziko & Dopad & Zmiernenie \\
\midrule
Obmedzený prístup k dátam & oneskorenie WP2 & verejné datasety, syntetické dáta, včasné rokovanie s partnermi \\
Nízka kvalita anotácie & skreslené výsledky & guidelines, viac anotátorov, kappa, kalibrácia \\
Nezvyknutý nepomer tried & nízky recall & cost-sensitive learning, resampling, kalibrácia \\
Posun distribúcie a LLM drift & pokles presnosti & časové testy, priebežné učenie, adversariálne testovanie \\
GDPR a citlivé údaje & právne riziko & minimalizácia, anonymizácia, súhlasy, etická kontrola \\
Odchod člena tímu & strata znalostí & dokumentácia, párové programovanie, zastupiteľnosť \\
Závislosť na partnerovi & oneskorenie pilotu & viac partnerov, offline pilot, jasné zmluvy \\
\bottomrule
\end{tabularx}
\end{table}

\subsection{Rozpočet}

Rozpočet vychádza z 36 mesiacov riešenia a je detailne rozpísaný vo formulári VV-C. Celkové náklady sú \budgetTotal\ EUR, z toho \budgetAPVV\ EUR z APVV a \budgetCoFin\ EUR zo spolufinancovania. Najväčšiu časť tvoria mzdové náklady a odvody, nasledujú služby, cestovné a nepriame náklady.

\begin{table}[h!]
\centering
\small
\begin{tabularx}{\textwidth}{@{}X r@{}}
\toprule
Kategória & Spolu (EUR) \\
\midrule
Mzdové náklady a ostatné osobné náklady & 140\,000 \\
Zdravotné a sociálne poistenie & 49\,000 \\
Cestovné náklady & 15\,000 \\
Materiál & 12\,000 \\
Služby & 30\,000 \\
Energie, vodné, stočné, komunikácie & 4\,000 \\
Náklady na popularizáciu & 8\,000 \\
Bežné nepriame náklady & 42\,000 \\
\midrule
Celkové náklady & 300\,000 \\
\bottomrule
\end{tabularx}
\end{table}

Rozpočet je primeraný cieľom: najviac prostriedkov smeruje na ľudskú prácu (anotácia, modelovanie, evaluácia), ktorá je pre projekt kritická. Cestovné náklady pokrývajú konferencie a stretnutia s partnermi, služby zahŕňajú cloudový výpočet a právne poradenstvo. Pred finálnym podaním rozpočet preveríme podľa taríf organizácie a podmienok výzvy.

\subsection{Infraštruktúra}

FIIT STU disponuje výpočtovou infraštruktúrou, laboratóriami a prístupom k vedeckým databázam. Pre projekt využijeme servery s GPU pre trénovanie transformerov, úložisko pre datasety a zabezpečené prostredie pre prácu s citlivými dátami. Personálnu infraštruktúru tvoria štyria študenti, zodpovedný riešiteľ a administratívna podpora fakulty. V prípade potreby doplníme cloudové zdroje z rozpočtu projektu.

{\small
\bibliographystyle{unsrt}
\bibliography{literatura}
}

\end{document}
